%----------------------------------------------------------------------------
\chapter{Archiving the code on Zenodo}\label{app:zenodo}
%----------------------------------------------------------------------------
A GitHub link can move, change or disappear; a reviewer checking \textbf{correctness} (criterion D/2) needs the exact code your results came from.
\href{https://zenodo.org/}{Zenodo}, run by CERN, archives a GitHub release permanently and gives it a DOI, so cite the archived release, not the repository.
These steps follow Zenodo's documentation as of September 2026; if the interface differs, the linked pages are authoritative.

\section{Prepare the repository}
\begin{enumerate}
  \item A \texttt{README.md} that says how to install the code and reproduce every measurement in \tplref{chap:evaluation}.
  \item A \texttt{LICENSE} file; ask your supervisor which license to use.
  \item A \texttt{.zenodo.json} file in the repository root (\cref{lst:zenodo-json}). Zenodo takes the record's title, authors, license and links from it (\href{https://help.zenodo.org/docs/github/describe-software/zenodo-json/}{format and example}). Check that it is valid JSON before releasing.
  \item If the repository has a \texttt{CITATION.cff}, keep its fields in line with \texttt{.zenodo.json}; Zenodo uses \texttt{.zenodo.json} when both exist.
\end{enumerate}

\begin{listing}[!ht]
\begin{minted}{json}
{
  "title": "TODO Artifact for: <thesis title>",
  "description": "Code and measurement scripts for the TDK paper <thesis title>.",
  "upload_type": "software",
  "version": "1.0.0",
  "publication_date": "2026-11-01",
  "license": "MIT",
  "creators": [
    { "name": "TODO Surname, Given name", "affiliation": "Budapest University of Technology and Economics", "orcid": "0000-0000-0000-0000" }
  ],
  "keywords": ["TODO"]
}
\end{minted}
\caption{A minimal \texttt{.zenodo.json}.}\label{lst:zenodo-json}
\end{listing}

\section{Publish the first release}
\begin{enumerate}
  \item Sign in to \href{https://zenodo.org/}{Zenodo} with the GitHub account that owns the repository. Open the profile menu, then \emph{GitHub}; click \emph{Sync now}, find the repository, and switch its toggle on (\href{https://help.zenodo.org/docs/github/enable-repository/}{enable a repository}).
  \item Commit the prepared files. On GitHub, open \emph{Releases}, then \emph{Draft a new release}: create a new tag such as \texttt{v1.0.0} on the commit your measurements used, add release notes, and publish.
  \item Back on Zenodo's GitHub page, open the repository and wait for processing. Open the new record and check its files and metadata; a failed release shows an \emph{Errors} section (\href{https://help.zenodo.org/docs/github/archive-software/github-upload/}{archive a release}).
\end{enumerate}
If the repository cannot be connected to Zenodo, upload a prepared archive instead: \emph{+}, then \emph{New upload}, fill in the metadata, \emph{Preview}, \emph{Publish} (\href{https://help.zenodo.org/docs/deposit/create-new-upload/}{upload instructions}).

\section{Cite it}
Each release gets its own \emph{version DOI}; the \emph{concept DOI} stands for all versions together (\href{https://support.zenodo.org/help/en-gb/1-upload-deposit/97-what-is-doi-versioning}{DOI versioning}).
\begin{itemize}
  \item In the paper, cite the \textbf{version DOI} of the release your results came from, as the \texttt{artifact} entry in \texttt{bib/references.bib} shows, and cite it in \tplref{chap:implementation}.
  \item In \texttt{README.md}, put the badge of the \textbf{concept DOI}, so it always leads to the newest version. The record page's DOI badge dialog gives the Markdown.
  \item Copy both DOIs from the Zenodo record; never guess a DOI by changing its number.
\end{itemize}

\section{Release a new version}
Fix the code, update \texttt{version} and \texttt{publication\_date} in \texttt{.zenodo.json}, commit, and publish a new GitHub release with a new tag (for example \texttt{v1.0.1}). Zenodo adds it to the same record as a new version with a new version DOI (\href{https://help.zenodo.org/docs/deposit/manage-versions/}{manage versions}). Never move an old tag. If the paper's results now come from the new version, cite its version DOI.
